% !TEX program = xelatex
% ---------- 共享 preamble ----------
% slides 与 script 共用。仅拉 _figpreamble.tex 的颜色宏 + tikz + pgfplots + subcaption。
% 不 \input 毕业论文/preamble.tex（避免 ctexbook 与 ctexbeamer/ctexart 冲突）。

% ---------- 字体 ----------
\def\ProjectFontDir{../fonts}
% Project-local font setup for XeLaTeX.
% Callers may set \ProjectFontDir before \input'ing this file.
\providecommand{\ProjectFontDir}{../fonts}

\setmainfont[
  Path = \ProjectFontDir/tex-gyre/,
  Extension = .otf,
  UprightFont = *-regular,
  BoldFont = *-bold,
  ItalicFont = *-italic,
  BoldItalicFont = *-bolditalic,
  Ligatures = TeX
]{texgyretermes}

\setCJKmainfont[
  Path = \ProjectFontDir/fandol/,
  Extension = .otf,
  BoldFont = FandolSong-Bold,
  ItalicFont = FandolKai-Regular
]{FandolSong-Regular}
\setCJKsansfont[
  Path = \ProjectFontDir/fandol/,
  Extension = .otf,
  BoldFont = FandolHei-Bold,
  ItalicFont = FandolKai-Regular
]{FandolHei-Regular}
\setCJKmonofont[
  Path = \ProjectFontDir/fandol/,
  Extension = .otf,
  ItalicFont = FandolKai-Regular
]{FandolFang-Regular}

\setCJKfamilyfont{zhsong}[
  Path = \ProjectFontDir/fandol/,
  Extension = .otf,
  BoldFont = FandolSong-Bold
]{FandolSong-Regular}
\setCJKfamilyfont{zhsongb}[
  Path = \ProjectFontDir/fandol/,
  Extension = .otf
]{FandolSong-Bold}
\setCJKfamilyfont{zhhei}[
  Path = \ProjectFontDir/fandol/,
  Extension = .otf,
  BoldFont = FandolHei-Bold,
  ItalicFont = FandolKai-Regular
]{FandolHei-Regular}
\setCJKfamilyfont{zhheib}[
  Path = \ProjectFontDir/fandol/,
  Extension = .otf
]{FandolHei-Bold}
\setCJKfamilyfont{zhfs}[
  Path = \ProjectFontDir/fandol/,
  Extension = .otf
]{FandolFang-Regular}
\setCJKfamilyfont{zhkai}[
  Path = \ProjectFontDir/fandol/,
  Extension = .otf
]{FandolKai-Regular}

\ExplSyntaxOn
\cs_if_exist:NT \ctex_punct_set:n
  {
    \ctex_punct_set:n { fandol }
    \ctex_punct_map_family:nn   { \CJKrmdefault         } { zhsong  }
    \ctex_punct_map_family:nn   { \CJKsfdefault         } { zhhei   }
    \ctex_punct_map_family:nn   { \CJKttdefault         } { zhfs    }
    \ctex_punct_map_bfseries:nn { \CJKrmdefault, zhsong } { zhsongb }
    \ctex_punct_map_bfseries:nn { \CJKsfdefault, zhhei  } { zhheib  }
    \ctex_punct_map_itshape:nn  { \CJKrmdefault         } { zhkai   }
  }
\ExplSyntaxOff

\providecommand{\songti}{}
\providecommand{\heiti}{}
\providecommand{\fangsong}{}
\providecommand{\kaishu}{}
\renewcommand{\songti}{\CJKfamily{zhsong}}
\renewcommand{\heiti}{\CJKfamily{zhhei}}
\renewcommand{\fangsong}{\CJKfamily{zhfs}}
\renewcommand{\kaishu}{\CJKfamily{zhkai}}


% ---------- 宏包 ----------
\usepackage{amsmath,amssymb}
\usepackage{booktabs}
\usepackage{array}
\usepackage{multirow}
\usepackage{tikz}
\usetikzlibrary{shapes.geometric, arrows.meta, positioning, fit, backgrounds, calc, patterns, decorations.pathreplacing}
\usetikzlibrary{external}
\tikzexternalize
\usepackage{pgfplots}
\pgfplotsset{compat=1.18}
\usepgfplotslibrary{fillbetween}
\pgfplotsset{table/search path={../图片/, ../图片/data/, ../}}
\usepackage{xcolor}
\usepackage{colortbl}
\usepackage{caption}
\usepackage{subcaption}
\usepackage{adjustbox}
\usepackage[absolute,overlay]{textpos}
\usepackage[skins]{tcolorbox}

% ---------- 统一调色板（与 图片/_figpreamble.tex 同步） ----------
\definecolor{deepblue}{RGB}{81, 132, 178}
\definecolor{lightblue}{RGB}{170, 212, 248}
\definecolor{skyblue}{RGB}{154, 220, 255}
\definecolor{inputyellow}{RGB}{255, 248, 154}
\definecolor{outputpink}{RGB}{255, 178, 166}
\definecolor{improvedpink}{RGB}{255, 138, 174}
\definecolor{lightpink}{RGB}{241, 167, 181}
\definecolor{deeppink}{RGB}{213, 82, 118}
\definecolor{bglight}{RGB}{242, 245, 250}
\definecolor{lightgreen}{RGB}{180, 230, 180}
\definecolor{lightorange}{RGB}{255, 210, 140}
\definecolor{dangerred}{RGB}{230, 100, 100}
\definecolor{vibrantgreen}{RGB}{46, 175, 80}
\definecolor{vibrantorange}{RGB}{240, 130, 30}
\definecolor{vibrantred}{RGB}{220, 50, 50}

% ---------- 答辩页浮动信息卡 ----------
\tcbset{
  slideMetricCard/.style={
    enhanced,
    colback=white,
    colframe=deepblue!70!black,
    coltitle=white,
    colbacktitle=deepblue!85!black,
    fonttitle=\bfseries\scriptsize,
    boxrule=0.5pt,
    arc=2pt,
    left=4pt,
    right=4pt,
    top=3pt,
    bottom=0pt,
    toptitle=1pt,
    bottomtitle=1pt,
    titlerule=0pt,
    drop fuzzy shadow=black!18
  }
}

% ---------- 统一 tikz 样式（与 图片/_figpreamble.tex 同步） ----------
\tikzset{
  mod/.style            = {rectangle, rounded corners=2pt, draw=deepblue, fill=lightblue, minimum height=0.9cm, align=center, font=\small},
  modimp/.style         = {rectangle, rounded corners=2pt, draw=deeppink, fill=improvedpink, minimum height=0.9cm, align=center, font=\small},
  modin/.style          = {rectangle, draw=deepblue, fill=inputyellow, minimum height=0.9cm, align=center, font=\small},
  modout/.style         = {rectangle, draw=deeppink, fill=outputpink, minimum height=0.9cm, align=center, font=\small},
  axx/.style            = {->, thick, black!80},
  flow/.style           = {-Stealth, thick, deepblue},
  flowimp/.style        = {-Stealth, thick, deeppink},
  obsstatic/.style      = {fill=dangerred!70, draw=dangerred, thick},
  obsdyn/.style         = {fill=lightorange, draw=dangerred, thick},
  obsbuf/.style         = {draw=dangerred!70, dashed, thick},
  bandok/.style         = {fill=lightgreen, opacity=0.55},
  bandbad/.style        = {fill=dangerred!30, opacity=0.55},
  egopt/.style          = {fill=deepblue, draw=deepblue},
  reflinestyle/.style   = {deepblue, very thick},
}
% ---------- 论文实验宏（与 毕业论文/ 同步） ----------
\InputIfFileExists{../毕业论文/_experiment_metrics.tex}{}{}
\InputIfFileExists{../毕业论文/_experiment_context.tex}{}{}
\InputIfFileExists{../毕业论文/_ablation_macros.tex}{}{}
\InputIfFileExists{../毕业论文/_sensitivity_macros.tex}{}{}

% ---------- 算法常量宏（slides 层级，非自动生成） ----------
\newcommand{\GroupKmax}{15}
\newcommand{\GroupCThreeK}{18}
\newcommand{\GroupCThreeCombCount}{26}
\newcommand{\ComplexityOldBase}{$O(2^k)$}
\newcommand{\ComplexityNewBase}{$O(k\log k)$}
